% Lay summary --- required by the Registrar since 2026.
% "Addressing a lay audience, provide a 250-word maximum description of the
%  significance of your dissertation."
% DRAFT, derived from abstract.tex -- please read and make it your own.
\begin{laysummary}

Before a computer can act on what we mean, we have to put it into words: type a
prompt, dictate a message, speak to an assistant. Words are a narrow channel.
We speak far more slowly than we think, and much of what we know---a movement,
an image, an intention---resists being said at all. This dissertation asks what
becomes possible if a machine can read intent more directly, from brain
activity itself.

The work proceeds in three steps. First, it builds machine-learning models that
respect how the brain is actually organized, rather than borrowing designs
meant for photographs and text. Some of these models discover structure that
existing brain atlases do not record; others build known biological structure,
such as the mirror symmetry between imagined left- and right-hand movement,
directly into the architecture.

Second, it connects recorded brain activity to meaning. Brain data are scarce
and expensive, so rather than collecting more, this work reuses a decade of
existing image descriptions, linked to brain recordings through the pictures
people were shown while being scanned. A model trained this way answers
open-ended questions about what a person saw or remembered, and reconstructs
the scene they were looking at.

Third, it turns those representations into action, learning to control a robot
from continuously recorded brain signals and video, without hand-labeled
trials.

Together these results sketch a path toward assistive technology for people who
cannot easily speak or move, and toward systems that respond to intention
rather than instruction.

\end{laysummary}
